%%%%%%%%%%%%%%%%%%%%%%%%%%%%%%%%%%%%%%%%%%%%%%%%%%%%%%%%%%%%%%%%%%%%%%%%%%%%%%%%
%   Copyright 2009, Oracle and/or its affiliates.
%   All rights reserved.
%
%
%   Use is subject to license terms.
%
%   This distribution may include materials developed by third parties.
%
%%%%%%%%%%%%%%%%%%%%%%%%%%%%%%%%%%%%%%%%%%%%%%%%%%%%%%%%%%%%%%%%%%%%%%%%%%%%%%%%

\section{\acffdcore}\label{acffd-calculus}

In this section, we define a Fortress core calculus with
acyclic type hierarchy and field definitions inside object definitions.
We call this calculus \emph{\acffdcore}.
\acffdcore\ is an extension of \basiccore\ with 
acyclic type hierarchy and field definitions inside object definitions.
\revision{revival-calculi}{This calculus predates instantiation exclusion
(\secref{trait-types}, \secref{trait-decls}).
Conditions (2) and (3) of its rule \acyclicRule\ admit one
self-extension, whose static arguments may be the parameters' own bounds,
so that a trait that extends an instantiation of itself at its parameter's
bound, such as
\EXP{\KWD{trait}\:\TYP{T}\llbracket{}X \KWD{extends} \TYP{Object}\rrbracket
\KWD{extends} \TYP{T}\llbracket\TYP{Object}\rrbracket \KWD{end}},
is well formed, and \EXP{\TYP{T}\llbracket{}A\rrbracket} is then a subtype
of \EXP{\TYP{T}\llbracket\TYP{Object}\rrbracket} for every \VAR{A};
Fortress no longer allows such a trait, for the reason the covariance
example of \secref{where-clauses} is not allowed.
The calculus, its rules and its soundness claim are as the team wrote and
proved them, and are not changed.}

\input{\macros/acffd-calculus-macros}
\input{\calculi/acffd/syntax}
\input{\calculi/acffd/dynamic}
\input{\calculi/acffd/static}
